\documentclass[11pt]{article}
\usepackage[a4paper,margin=1in]{geometry}
\usepackage{amsmath,amssymb,amsthm,microtype}
\newtheorem{theorem}{Theorem}[section]
\theoremstyle{remark}\newtheorem{remark}[theorem]{Remark}
\newcommand{\Tr}{\operatorname{Tr}}
\title{Annular Tail Reduction to the Moving Spectral Head}
\author{Bin Wang}\date{July 2026}
\begin{document}\maketitle
\begin{abstract}
The direct annular criterion of RH-300 was left inactive because no norm
estimate was known for the actual complement-to-anchor logarithmic mismatch.
We prove that its high-order part is already controlled.  For every fixed
$1.4<\rho<\rho_*=1.426787\ldots$, the noisy modulus-complement tail and the
deterministic target tail beyond
$m_\sigma=\lceil4\log(1/\sigma)\rceil$ vanish in both $H^\infty(\rho)$ and
$H^2(\rho)$ with explicit bounds.  Consequently full annular convergence is
equivalent to convergence of the moving polynomial head below $m_\sigma$.
This locates the remaining obstruction without proving that the moving head
converges.
\end{abstract}

\section{Setup}
Put
\[
 q=\frac12,\qquad R=\frac75,\qquad
 q_*=(0.85\lambda)^{-1}=0.700875225854775\ldots,
 \qquad \rho_*=q_*^{-1}.
\]
Let $C_\sigma$ be the modulus-complete spectral complement from RH-282,
with squared spectral mass $M_\sigma\le\sigma^{-1}$, and write
\[
 \tau_{\sigma,n}=\Tr C_\sigma^n,
 \qquad |\tau_{\sigma,n}|\le\sigma^{-1}q^{n-2}.
\]
The deterministic numerator anchors satisfy the all-order RH-267 envelope
$|a_n|<48q_*^n$.  Define
\[
 g_\sigma(z)=\sum_{n\ge2}
 \frac{\tau_{\sigma,n}-a_n}{n}z^n,
 \qquad
 g_\sigma^{<m}(z)=\sum_{2\le n<m}
 \frac{\tau_{\sigma,n}-a_n}{n}z^n.
\]
For $q\rho<1$ and $q_*\rho<1$, both series are well defined on the closed
disk of radius $\rho$.

\section{Tail theorem}
\begin{theorem}[slope-four annular tail reduction]\label{thm:tail}
Fix $R<\rho<\rho_*$ and put
$m_\sigma=\lceil4\log(1/\sigma)\rceil$,
$x=q\rho$, and $y=q_*\rho$.  Then
\begin{align*}
 \|g_\sigma-g_\sigma^{<m_\sigma}\|_{H^\infty(\rho)}
 &\le
 \frac{4\sigma^{-1}x^{m_\sigma}}
 {m_\sigma(1-x)}
 +\frac{48y^{m_\sigma}}
 {m_\sigma(1-y)},\\
 \|g_\sigma-g_\sigma^{<m_\sigma}\|_{H^2(\rho)}
 &\le
 \frac{4\sigma^{-1}x^{m_\sigma}}
 {m_\sigma\sqrt{1-x^2}}
 +\frac{48y^{m_\sigma}}
 {m_\sigma\sqrt{1-y^2}}.
\end{align*}
Both right sides tend to zero.  More precisely, apart from the common
$1/\log(1/\sigma)$ factor, their two powers of $\sigma$ are
\[
 4\log(2/\rho)-1
 \quad\hbox{and}\quad
 4\log(\rho_*/\rho).
\]
Hence for $X=H^\infty(\rho)$ or $H^2(\rho)$,
\[
 \|g_\sigma\|_X\longrightarrow0
 \quad\Longleftrightarrow\quad
 \|g_\sigma^{<m_\sigma}\|_X\longrightarrow0.
\]
\end{theorem}
\begin{proof}
For the noisy tail, use
$|\tau_{\sigma,n}|\le4\sigma^{-1}q^n$ and
\[
 \sum_{n\ge m}\frac{x^n}{n}
 \le\frac{x^m}{m(1-x)},\qquad
 \left(\sum_{n\ge m}\frac{x^{2n}}{n^2}\right)^{1/2}
 \le\frac{x^m}{m\sqrt{1-x^2}}.
\]
The same estimates with $48y^n$ give the target terms.  Since
$m_\sigma\ge4\log(1/\sigma)$ and $x,y<1$, the displayed powers follow.
They are positive throughout $R<\rho<\rho_*$.  The equivalence is the
triangle inequality applied to the vanishing tail.
\end{proof}

\section{Numerical protocol and scope}
The reproducible computation evaluates the two certified bounds at
$\rho=1.405,1.41,1.42$.  At $\rho=1.41$ the noisy and target exponents are
$0.3982299047\ldots$ and $0.0473427916\ldots$.  The small target exponent
explains why finite-noise upper bounds can remain large even though the
asymptotic theorem is strict.

\begin{remark}
The theorem closes only the analytic high-order tail.  It supplies no
estimate for the moving polynomial $g_\sigma^{<m_\sigma}$ and therefore does
not activate RH-300.  Gates A--E remain false/open; no Hilbert--Polya,
Riemann-zero, zeta-divisor, von Mangoldt trace, or RH conclusion is made.
\end{remark}
\end{document}
